\section{Experimentos}
\label{sec:exp}

Este capítulo presenta la evaluación experimental en tres bloques: la
comparación controlada de las tres arquitecturas
(Sección~\ref{sec:arch-comp}); seis casos de prueba de complejidad creciente
---del régimen 1D estacionario a datos reales de laboratorio--- que conforman
el banco de pruebas (Secciones~\ref{sec:exp1} a~\ref{sec:exp6}); y la ablación
de la forma del residuo SWE (Sección~\ref{sec:ablation}). Cada caso se
especifica a continuación; sus resultados se completarán con las corridas del
estudio comparativo de arquitecturas.

\subsection{Comparación de arquitecturas}
\label{sec:arch-comp}

Este estudio compara los tres diseños de arquitectura de la
Sección~\ref{sec:arquitecturas} --- la propuesta de dos redes (A1), la
monolítica estilo Dazzi (A2) y la de una red por campo estilo Ohara (A3) ---
mediante un
\emph{estudio de escalamiento}: cada diseño se instancia a tres presupuestos
de parámetros (pequeño $\approx$20K, medio $\approx$100K, grande
$\approx$500K) y se evalúa en tres casos de complejidad creciente --- el bump
subcrítico 1D estacionario (Exp.~1), la cuenca de Thacker 1D transitoria
(Exp.~2) y los dos cilindros 2D (Exp.~3) ---, bajo el protocolo de comparación
justa de la Sección~\ref{sec:protocolo}. El resultado es doble: (i) a
presupuesto de parámetros igualado, qué diseño recupera mejor la batimetría
--- comparación que aísla el diseño de la capacidad ---; y (ii) la curva de
precisión ($\RMSE$, $\NRMSE$, $R^2$ de $\zb$) frente al número de parámetros
de cada diseño, que sitúa a la propuesta en la frontera precisión-costo. Se
reporta además el costo de cómputo (tiempo de entrenamiento y VRAM pico). Los
tres diseños se ordenan por su \emph{grado de separación de redes}: monolítica
(A2), flujo agrupado con batimetría aparte (A1) y una red por cada campo (A3).
El objetivo es ubicar a la propuesta en ese eje --- si darle red propia a
$\zb$ y agrupar el flujo con embedding de Fourier le da ventaja sobre la red
monolítica y sobre la separación total por campo, y en qué rango de
presupuestos.

\paragraph{Resultados.}
\pendiente{tabla del estudio de escalamiento (3 diseños $\times$ 3 presupuestos
$\times$ 3 casos) y figura de la frontera precisión--parámetros.}

% ──────────────────────────────────────────────────────────────
\subsection{Exp.\ 1 --- Bump subcrítico 1D (estacionario)}
\label{sec:exp1}

\paragraph{Setup.}
Bump parabólico de 200\,mm en dominio $[-10, 10]$\,m, descarga
$q = 4.42$\,m$^2$/s, profundidad aguas abajo $h_{\text{down}} = 2.0$\,m,
sin fricción. Rango de Froude: $[0.50, 0.63]$. Depresión superficial:
93\,mm. El planteamiento y la solución analítica (ecuación de Bernoulli) se
detallan en el Apéndice~\ref{app:exp1}.

\paragraph{Resultados.}
\pendiente{resultados del estudio de arquitecturas.}

% ──────────────────────────────────────────────────────────────
\subsection{Exp.\ 2 --- Cuenca de Thacker 1D (transitorio)}
\label{sec:exp2}

\paragraph{Setup.}
Cuenca parabólica $\zb = 0.5(x^2 - 1)$ en $[-2, 2]$\,m, profundidad
central $h_0 = 0.5$\,m, período de oscilación $T = 2.006$\,s.
Problema cerrado sin condiciones de borde de flujo. El planteamiento y la
solución analítica de Thacker se detallan en el Apéndice~\ref{app:exp2}.

\paragraph{Resultados.}
\pendiente{resultados del estudio de arquitecturas.}

% ──────────────────────────────────────────────────────────────
\subsection{Exp.\ 3 --- Dos cilindros 2D (transitorio)}
\label{sec:exp3}

\paragraph{Setup.}
Dos obstáculos cilíndricos en un dominio 2D; observaciones $\eta$, $u$ y $v$
sobre una malla $40\times40$ y 15 instantáneas temporales. El planteamiento y
la solución de referencia numérica se detallan en el Apéndice~\ref{app:exp3}.

\paragraph{Resultados.}
\pendiente{resultados del estudio de arquitecturas.}

% ──────────────────────────────────────────────────────────────
\subsection{Exp.\ 4 --- Topografía variable 2D (Tian dT10)}
\label{sec:exp4}

\paragraph{Setup.}
Réplica fiel del caso dinámico \emph{dT10} de Tian et al.\ \cite{tian2025}
(Sección~3.3): masa de agua sobre topografía oscilatoria 2D que parte de un
estado desbalanceado ($h(x,y,0) = \zb$) y se relaja bajo gravedad con bordes
periódicos, sin forzamiento externo. El planteamiento y la solución de
referencia numérica se detallan en el Apéndice~\ref{app:exp4}.

\paragraph{Resultados.}
\pendiente{resultados del estudio de arquitecturas.}

% ──────────────────────────────────────────────────────────────
\subsection{Exp.\ 5 --- Cuenca cerrada (paraboloide de Thacker)}
\label{sec:exp5}

\paragraph{Setup.}
Cuenca cuyo borde emerge sobre el nivel medio (75.5\% del dominio siempre
seco) y \emph{sin} forzamiento de borde externo. El planteamiento y la
solución analítica de Thacker 2D se detallan en el Apéndice~\ref{app:exp5}.

\paragraph{Resultados.}
\pendiente{resultados del estudio de arquitecturas.}

% ──────────────────────────────────────────────────────────────
\subsection{Exp.\ 6 --- Datos reales (canal de oleaje de Angel)}
\label{sec:exp6}

\paragraph{Setup.}
Datos \emph{reales} del canal de Hamburgo de Angel et al.\ \cite{angel2024},
ventana de oleaje establecido $t\in[40,60]$\,s; el sensor S1 ($x=1.5$\,m)
impone la entrada como Dirichlet \emph{soft} (penalización MSE con el
mismo peso que la observación interior, no es una BC dura --- el método
adjoint de Angel sí impone la BC dura a través del solver directo, lo
que parcialmente explica el cierre del gap NRMSE) y S2 ($x=3.5$\,m),
fuera del bump, es la observación interior; PINN transitorio 1D con
drag de fondo lineal y $\zb\ge0$.

\paragraph{Resultados.}
\pendiente{resultados del estudio de arquitecturas.}

% ──────────────────────────────────────────────────────────────
\subsection{Estudio de ablación: forma de ecuación SWE}
\label{sec:ablation}

Se comparan las tres formas del residuo SWE (Sección~\ref{sec:method})
en el caso del bump subcrítico (Exp.\ 1), con la misma arquitectura,
pesos de pérdida y presupuesto de entrenamiento. Cada forma se entrena
con semillas aleatorias independientes para evaluar la robustez.

\paragraph{Resultados.}
\pendiente{tabla comparativa de las tres formas del residuo (RMSE, $R^2$ y
robustez entre semillas).}
